\documentclass[10pt]{article}

% ---------- Document Identity (update these only) ----------
\newcommand{\DocStage}{}
\newcommand{\DocTitle}{Your Road to Tier One SOF}
\newcommand{\ProgramName}{182-Week Civilian-to-Selection Readiness Pipeline}
\newcommand{\ProgramSubtitle}{}

% ---------- Page ----------
\usepackage[legalpaper,left=0.5in,right=0.5in,top=0.6in,bottom=0.45in]{geometry}

% ---------- Typography ----------
\usepackage[T1]{fontenc}
\usepackage[utf8]{inputenc}
\usepackage{libertinus}
\usepackage{microtype}
\usepackage{parskip}
\usepackage{ragged2e}
\usepackage{textcomp}
\AtBeginDocument{\color{black}}
\AtBeginDocument{\pagecolor{white}}
\AtBeginDocument{\justifying}

% ---------- Landscape pages ----------
\usepackage{pdflscape}

% ---------- Color ----------
\usepackage[table]{xcolor}
\definecolor{StageBlue}{HTML}{1F4E79}
\definecolor{StageBlueDark}{HTML}{163A5A}
\definecolor{LightGray}{HTML}{F2F4F7}
\definecolor{MidGray}{HTML}{D9DEE5}
\definecolor{RuleGray}{HTML}{C7CED8}
\definecolor{HeaderInk}{HTML}{000000}
\definecolor{Ink}{HTML}{000000}

% ---------- Utility ----------
\usepackage{enumitem}
\sloppy
\setlength{\emergencystretch}{2em}

% ---------- Boxes / UI ----------
\usepackage[most]{tcolorbox}

\tcbset{
  charterTitle/.style={
    enhanced,
    colback=StageBlue!4,
    colframe=StageBlue,
    boxrule=1.25pt,
    arc=3pt,
    left=16pt,right=16pt,top=14pt,bottom=14pt,
    borderline north={3.2pt}{0pt}{StageBlue},
    borderline west={4pt}{0pt}{StageBlue},
    borderline south={1.25pt}{0pt}{StageBlue},
    drop shadow={black!25!white},
  },
  charterRule/.style={
    enhanced,
    colback=LightGray,
    colframe=RuleGray,
    boxrule=0.85pt,
    arc=3pt,
    left=12pt,right=12pt,top=10pt,bottom=10pt,
    borderline west={3pt}{0pt}{StageBlue},
  },
  charterBadge/.style={
    enhanced,
    colback=StageBlue,
    coltext=white,
    colframe=StageBlueDark,
    boxrule=0.6pt,
    arc=2pt,
    left=6pt,right=6pt,top=3pt,bottom=3pt,
  }
}
\newtcbox{\Badge}{nobeforeafter,charterBadge}

% ---------- Lists ----------
\setlist[itemize]{leftmargin=1.5em, itemsep=0.25em, topsep=0.25em}
\setlist[enumerate]{leftmargin=1.8em, itemsep=0.25em, topsep=0.25em}

% ---------- TOC ----------
\usepackage{tocloft}
\setcounter{tocdepth}{3}
\setlength{\cftbeforesecskip}{3pt}
\setlength{\cftbeforesubsecskip}{2pt}
\setlength{\cftbeforesubsubsecskip}{0pt}
\renewcommand{\cftsecfont}{\bfseries}
\renewcommand{\cftsecpagefont}{\bfseries}
\renewcommand{\cftsubsecfont}{\normalfont}
\renewcommand{\cftsubsecpagefont}{\normalfont}
\renewcommand{\cftsubsubsecfont}{\normalfont}
\renewcommand{\cftsubsubsecpagefont}{\normalfont}
\setlength{\cftsecindent}{0pt}
\setlength{\cftsubsecindent}{1.25em}
\setlength{\cftsubsubsecindent}{2.8em}
\setlength{\cftsecnumwidth}{2.1em}
\setlength{\cftsubsecnumwidth}{2.8em}
\setlength{\cftsubsubsecnumwidth}{3.6em}

% Remove the built-in TOC heading so only the custom heading is shown
\makeatletter
\renewcommand{\tableofcontents}{\@starttoc{toc}}
\makeatother

% ---------- Header/Footer: page number only ----------
\usepackage{fancyhdr}
\pagestyle{fancy}
\fancyhf{}
\renewcommand{\headrulewidth}{0pt}
\renewcommand{\footrulewidth}{0pt}
\fancyfoot[C]{\thepage}

% ---------- Section formatting ----------
\usepackage{titlesec}

\titleformat{\section}
  {\Large\bfseries\color{Ink}}
  {\thesection}{0.6em}{}
  [\vspace{0.25em}\textcolor{StageBlue}{\rule{\linewidth}{0.8pt}}]

\titleformat{\subsection}
  {\large\bfseries\color{Ink}}
  {\thesubsection}{0.6em}{}

\titleformat{\subsubsection}
  {\normalsize\bfseries\color{Ink}}
  {\thesubsubsection}{0.6em}{}

\titlespacing*{\section}{0pt}{0.95em}{0.35em}
\titlespacing*{\subsection}{0pt}{0.65em}{0.25em}
\titlespacing*{\subsubsection}{0pt}{0.55em}{0.2em}

% ---------- Table engine ----------
\usepackage{tabularray}
\UseTblrLibrary{booktabs}

% ---------- tabularray: GLOBAL table treatment ----------
\SetTblrInner{
  cells = {fg=black, bg=white, valign=t},
}

% ---------- Header helpers ----------
\newcommand{\Hdr}[1]{\shortstack[c]{\strut #1\strut}}

% ---------- Lines ----------
\newcommand{\TitleLine}{\textcolor{StageBlue}{\rule{0.98\linewidth}{0.95pt}}}
\newcommand{\SubTitleLine}{\textcolor{RuleGray}{\rule{0.98\linewidth}{0.65pt}}}

% ---------- Continued on next page ----------
\newcommand{\ContinuedOnNextPageInline}{%
  \par
  \vfill
  \noindent\textcolor{RuleGray}{\rule{\linewidth}{1pt}}\vspace{-2pt}
  \noindent\hfill\textit{Continued on next page}\par\vspace{-10pt}
  \noindent\textcolor{RuleGray}{\rule{\linewidth}{1pt}}
  \clearpage
}

% ---------- PDF hyperlinks ----------
\usepackage[hidelinks]{hyperref}
\hypersetup{
  colorlinks=false,
  pdfborder={0 0 0},
  linktoc=all,
  pdftitle={\DocTitle},
  pdfauthor={\ProgramName}
}

% =========================================================
\begin{document}

\subsection*{Phase 06 Card Header}
\phantomsection
\addcontentsline{toc}{subsection}{Phase 06 Card Header}

\begin{itemize}
\item Phase \textbf{ID}: Phase 06
\item Phase Name: Lactate Threshold and Power-Endurance Development
\item Duration weeks: 12
\item Version: v1.0
\item Stage Alignment: Stage 5
\item Governing Status: Draft — Non-Deployable
\end{itemize}

\subsection*{1. Phase Mission}
\phantomsection

Increase selection-relevant power-endurance and lactate-threshold expression while preserving durability and evidence validity by concentrating gate-grade tests into protected deload and test windows and enforcing anti-stacking posture across \textbf{RUN}, \textbf{RUCK}, and \textbf{SWIM} domains.

\subsection*{2. Phase Purpose}
\phantomsection

\begin{itemize}
\item Establish a controlled shift from “aerobic base emphasis” toward more consequential power-endurance verification while keeping validity and safety dominant.
\item Preserve Phase-to-phase continuity by maintaining \textbf{CAL} standards and durability controls while introducing or sustaining higher-cost locomotion verification (\textbf{RUN} and \textbf{RUCK}) under strict reslot discipline.
\item Set up the next phase by proving repeatable gate-grade outcomes under stable environments/tools and stable recovery posture (no novelty near gates).
\end{itemize}

\subsection*{3. Entry Criteria}
\phantomsection

\begin{itemize}
\item Entry requires admissible end-of-prior-phase evidence consistent with Stage 1.F (\textbf{VALID} sessions and \textbf{VALID} tests) and completion of the prior phase exit posture.
\item Default posture if evidence is incomplete, invalid, or not replicated: \textbf{HOLD} and reslot the missing or invalid gate tests to the next protected window per Stage 3.E; do not proceed on assumed readiness.
\end{itemize}

\subsection*{4. Operating Constraints}
\phantomsection

\begin{itemize}
\item Six sessions per week is mandatory. Risk is controlled by downshifting dose, modality, and intensity, not by collapsing frequency.
\item Required session elements must occur in this order for a session to be progression-eligible:
  \begin{itemize}
  \item Safety self-check first
  \item Warm-up
  \item Mobility
  \item Main work
  \item Cardio
  \item Cooldown
  \item Recovery notes and any required post-session notes per Stage 1.F
  \end{itemize}
\item Cardio minimum standard is minimum 30 minutes continuous per session.
\item 10,000 steps per day is a global requirement and must be logged.
\item Gate decisions use admissible evidence only. \textbf{INVALID} tests provide no gate credit and cannot support \textbf{ADVANCE}.
\item No underwater or breath-hold exposure unless governance prerequisites are satisfied; underwater is not a Phase 06 gate driver by default and must never be scheduled if governance cannot be guaranteed.
\end{itemize}

\subsection*{5. Primary Adaptations and Physiological Focus}
\phantomsection

\begin{itemize}
\item Increase power-endurance capacity while maintaining stable mechanics and recovery.
\item Improve pace/effort sustainability under controlled fatigue without converting verification into breathless risk behavior.
\item Preserve durability (feet, gait mechanics, pain trend stability) so higher-cost \textbf{RUN}/\textbf{RUCK} gates remain admissible.
\item Maintain \textbf{CAL} standards as a stable readiness anchor while locomotion gates increase consequence.
\item Strength proxy remains supporting only; Phase 06 does not use maximal strength testing.
\end{itemize}

\subsection*{6. Primary Modalities}
\phantomsection

\begin{itemize}
\item Emphasized domains: \textbf{RUN}; \textbf{RUCK}; \textbf{CAL}; \textbf{RECOVERY}; \textbf{DURABILITY}.
\item Supporting domains: \textbf{SWIM} (when pool access/verification is stable); \textbf{STR} (submax proxy only); \textbf{BODYCOMP} (trend stability).
\item De-emphasized / prohibited: maximal strength testing; unscheduled time trials in build weeks; any non-Stage 2 test constructs for gate claims; any unsupervised underwater/breath-hold exposure.
\end{itemize}

\clearpage
\begin{landscape}

\subsection*{7. Weekly Structure}
\phantomsection

\begingroup
\scriptsize
\setlength{\tabcolsep}{2pt}
\renewcommand{\arraystretch}{1.12}

\begin{longtblr}{
  width=\linewidth,
  rowhead=1,
  colspec = {
    Q[l,wd=0.70in]
    X[l,2.05]
    X[l,1.25]
    X[l,2.95]
    X[l,2.90]
  },
  hlines,
  vlines,
  cells = {hsep=6pt, vsep=6pt, halign=l, valign=t},
  cell{1-Z}{1-5} = {fg=black},
  row{odd} = {bg=white},
  row{even} = {bg=LightGray},
  row{1} = {bg=StageBlue!15, fg=black, font=\bfseries, halign=l, valign=t},
}
\Hdr{Session \#} &
\Hdr{Primary Focus} &
\Hdr{Main Work Domains} &
\Hdr{Cardio Modality/Intent} &
\Hdr{Notes/Risk Controls} \\

1 &
\textbf{RUN} technique and readiness emphasis &
\textbf{RUN}; \textbf{DURABILITY} &
Modality: step-based locomotion or treadmill run/walk per phase intent; Minutes: 30-45; RPE + Talk Test: controlled, non-breathless; Continuity: continuous planned &
Protect gait mechanics; no novelty footwear/route/tool near protected windows; downshift on pain drift. \\

2 &
\textbf{CAL} quality and trunk endurance continuity &
\textbf{CAL}; \textbf{DURABILITY} &
Modality: low-impact steady; Minutes: 30-45; RPE + Talk Test: easy-to-moderate; Continuity: continuous planned &
\textbf{CAL} battery standards must remain stable; artifact posture enforced when \textbf{CAL} is tested. \\

3 &
\textbf{RUCK} tolerance and foot interface governance &
\textbf{RUCK}; \textbf{DURABILITY} &
Modality: step-based locomotion (ruck exposure intent) or indoor substitute when unsafe; Minutes: 45-75; RPE + Talk Test: controlled; Continuity: continuous planned &
Foot hotspot monitoring is decisive; do not force load exposure when foot breakdown is present; reslot gates if compromised. \\

4 &
\textbf{RUN} power-endurance emphasis &
\textbf{RUN}; \textbf{RECOVERY} &
Modality: low-impact steady or controlled run/walk exposure; Minutes: 30-45; RPE + Talk Test: controlled short-sentence allowed only if safe; Continuity: continuous planned &
Avoid stacking high-cost \textbf{RUN} with high-cost \textbf{RUCK} within spacing posture; downshift if recovery markers degrade. \\

5 &
\textbf{SWIM} or \textbf{STR} supporting day (phase-dependent) &
\textbf{SWIM} or \textbf{STR}; \textbf{RECOVERY} &
Modality: pool swim when stable OR low-impact substitute when access/validity not possible; Minutes: 30-45; RPE + Talk Test: controlled; Continuity: continuous planned &
\textbf{SWIM} validity requires pool verification; if not stable, \textbf{SWIM} becomes non-gate-supporting or reslotted. \textbf{STR} is submax proxy only. \\

6 &
Consolidation and readiness stabilization &
\textbf{RECOVERY}; \textbf{DURABILITY} &
Modality: lowest-risk continuous Cardio; Minutes: 30-45; RPE + Talk Test: easy; Continuity: continuous planned &
Use to protect readiness approaching gates; no “make-up” behavior; preserve logs and stability. \\

\end{longtblr}

\endgroup

\subsection*{8. Progression Rules}
\phantomsection

\begin{itemize}
\item One progression lever at a time: do not increase locomotion dose and increase \textbf{CAL} density and increase \textbf{RUCK} consequence in the same week.
\item Preserve intent while reducing risk: regress ROM/leverage/speed/complexity before reducing session frequency.
\item Regression triggers: pain that alters mechanics, foot hotspot escalation, sleep collapse, unusual breathlessness, or repeated \textbf{MODIFIED}/\textbf{INVALID} sessions in the decision window → downshift and reslot gate attempts if validity would be compromised.
\item No novelty near protected windows: route/tool/footwear changes are treated as comparability breaks and may render evidence non-admissible until stabilized.
\end{itemize}

\subsection*{9. Deload and Test Placement}
\phantomsection

\begin{itemize}
\item Phase 06 uses Stage 3 protected windows. Week-in-phase placements for Phase 06 are: Week 1 (baseline), Week 4 (Deload+Test), Week 8 (Deload+Test), Week 12 (Deload+Test end gate).
\item Final placement and any reslot actions remain governed by Stage 3 integrated calendars and Stage 3.E adjustment doctrine.
\end{itemize}

\end{landscape}

\clearpage
\begin{landscape}

\subsection*{10. Test Battery}
\phantomsection

\begingroup
\scriptsize
\setlength{\tabcolsep}{2pt}
\renewcommand{\arraystretch}{1.12}

\begin{longtblr}{
  width=\linewidth,
  rowhead=1,
  colspec = {
    X[l,1.55]
    X[l,3.05]
    Q[l,wd=0.95in]
    X[l,3.35]
    Q[l,wd=1.35in]
  },
  hlines,
  vlines,
  cells = {hsep=6pt, vsep=6pt, halign=l, valign=t},
  cell{1-Z}{1-5} = {fg=black},
  row{odd} = {bg=white},
  row{even} = {bg=LightGray},
  row{1} = {bg=StageBlue!15, fg=black, font=\bfseries, halign=l, valign=t},
}
\Hdr{Test Timing} &
\Hdr{Test \textbf{ID}/Name} &
\Hdr{Domain} &
\Hdr{Validity Conditions} &
\Hdr{Pass Threshold Stage 2 Tier} \\

Week 1 baseline &
\textbf{C1} / \textbf{MET-2B-001} Bodyweight — Weekly Average (baseline capture event) &
\textbf{BODYCOMP} &
Standard conditions and method recorded; same-day logging &
\textbf{ENTRY} \\

Week 1 baseline &
\textbf{C7} / \textbf{MET-2B-013} Plank — Forearm &
\textbf{CAL} &
Standard form; timing method stable; same-day logging; if \textbf{CAL} artifact required by protocol posture, it must be compliant &
\textbf{ENTRY} \\

Week 4 mid-phase gate window &
\textbf{C8} / \textbf{MET-2B-014} 1.5-Mile Run Time Trial &
\textbf{RUN} &
Route or treadmill admissibility posture; machine/route identity logged; warm-up logged; no stop-rule violations &
\textbf{INTERMEDIATE} \\

Week 4 mid-phase gate window &
\textbf{C4} / \textbf{MET-2B-010} Push-Ups — 2-Minute Timed, Strict Standard &
\textbf{CAL} &
\textbf{CAL} artifact evidence required for \textbf{VALID} gate use; standardized conditions; warm-up logged &
\textbf{INTERMEDIATE} \\

Week 8 mid-phase gate window &
\textbf{C12} / \textbf{MET-2B-019} 6-Mile Ruck Time Trial — 35 lb dry &
\textbf{RUCK} &
Dry load verified; water logged separately; route/treadmill identity logged; warm-up logged; no stop-rule violations &
\textbf{INTERMEDIATE} \\

Week 8 mid-phase gate window &
\textbf{C13} / \textbf{MET-2B-023} Swim 500 m Time — No Fins &
\textbf{SWIM} &
Pool unit verified; no fins; allowable strokes; warm-up logged; rest rules adhered &
\textbf{INTERMEDIATE} \\

Week 12 end-phase gate window &
\textbf{C10} / \textbf{MET-2B-016} 3-Mile Run Time Trial &
\textbf{RUN} &
Route/treadmill identity logged; warm-up logged; no stop-rule violations; comparable conditions preserved &
\textbf{INTERMEDIATE} \\

Week 12 end-phase gate window &
\textbf{C12} / \textbf{MET-2B-020} 9-Mile Ruck Time Trial — 45 lb dry &
\textbf{RUCK} &
Dry load verified; water logged separately; route/treadmill identity logged; warm-up logged; no stop-rule violations &
\textbf{INTERMEDIATE} \\

Week 12 end-phase gate window &
\textbf{C6} / \textbf{MET-2B-012} Pull-Ups — Strict Dead-Hang &
\textbf{CAL} &
\textbf{CAL} artifact evidence required for \textbf{VALID} gate use; pronated strict standard; full ROM auditable &
\textbf{INTERMEDIATE} \\

\end{longtblr}

\endgroup

\end{landscape}
\clearpage

\subsection*{11. Exit Standards}
\phantomsection

\begin{itemize}
\item Phase 06 exit decision is evaluated at the end-phase protected window using admissible evidence only.
\item Intended tier posture alignment for Phase 06 exit: \textbf{INTERMEDIATE}.
\item Minimum posture: the end-phase gate tests listed in the Test Battery must meet \textbf{INTERMEDIATE} tier thresholds under \textbf{VALID} conditions, with replication posture governed by Stage 2 gate-use requirements.
\item If any gate-critical test is \textbf{INVALID}, it provides no gate credit and defaults to \textbf{HOLD} with reslot to the next protected window per Stage 3.E.
\item If durability constraints are active (gait-altering pain, active blister/hotspot progression, or unresolved mechanics drift), the correct posture is \textbf{HOLD} or \textbf{REGRESS} consistent with Stage 1.F.
\end{itemize}

\subsection*{12. Risk Controls and Stop Rules}
\phantomsection

\begin{itemize}
\item Stage 0.5 and Stage 1.F stop posture and escalation posture control this phase; this section adds phase-specific controls only.
\item \textbf{RUN} and \textbf{RUCK} are the highest-cost domains in Phase 06:
  \begin{itemize}
  \item enforce anti-stacking intent: do not schedule maximal \textbf{RUN} and maximal \textbf{RUCK} adjacent; reslot rather than compress.
  \item enforce environment safety: unsafe footing, severe heat, poor air quality triggers indoor substitution only when comparability can be preserved and logged; otherwise reslot.
  \end{itemize}
\item Foot integrity is a gate constraint:
  \begin{itemize}
  \item hotspots trending to blisters or active blisters constrain \textbf{RUCK} and \textbf{RUN} gate attempts; reslot the affected tests.
  \end{itemize}
\item \textbf{SWIM} is gate-grade only when pool verification and rest governance can be met; otherwise \textbf{SWIM} gate attempts are \textbf{INVALID} and must be reslotted.
\item Underwater remains governance-gated; do not schedule if governance cannot be guaranteed.
\end{itemize}

\clearpage
\begin{landscape}

\subsection*{13. Required Capabilities and Dependencies}
\phantomsection

\begingroup
\scriptsize
\setlength{\tabcolsep}{2pt}
\renewcommand{\arraystretch}{1.12}

\begin{longtblr}{
  width=\linewidth,
  rowhead=1,
  colspec = {
    X[l,1.85]
    X[l,3.15]
    Q[l,wd=1.25in]
    Q[l,wd=1.25in]
    X[l,3.30]
  },
  hlines,
  vlines,
  cells = {hsep=6pt, vsep=6pt, halign=l, valign=t},
  cell{1-Z}{1-5} = {fg=black},
  row{odd} = {bg=white},
  row{even} = {bg=LightGray},
  row{1} = {bg=StageBlue!15, fg=black, font=\bfseries, halign=l, valign=t},
}
\Hdr{Dependency \textbf{ID}} &
\Hdr{Required Capability Category and Tier} &
\Hdr{Due-By week-in-phase} &
\Hdr{Owner} &
\Hdr{Fail Action} \\

\textbf{DEP-1C-015} &
7.09 Health Monitoring Diagnostics and Biometrics — Tier 0 (minimum logging feasibility and safety awareness) &
Week 1 &
Medical &
\textbf{HOLD}. Do not proceed on unclear medical restriction posture; reslot gate attempts until cleared within governance boundaries. \\

\textbf{DEP-1C-007} \& \textbf{DEP-1C-008} &
7.11 Footwear Systems and Foot Care — Tier 1 (stable shoe/sock/lacing; hotspot controls) &
Week 1 &
Equipment Lead &
\textbf{HOLD} on \textbf{RUCK}/\textbf{RUN} gate attempts. Substitute low-impact Cardio modalities; reslot locomotion gates to next protected window. \\

\textbf{DEP-1C-006} &
7.18 Testing, Timing, Measurement, and Course-Setup Tools — Tier 1 (route/machine identity capture; timing method; load verification method) &
Week 1 &
Coach Lead &
Any gate attempt without stable setup is \textbf{INVALID} for gate use; reslot; maintain training-only exposures until setup is stable. \\

\textbf{DEP-1C-013} &
7.07 Load-Bearing Rucking and Carry Systems — Tier 1 (safe dry-load configuration; verification posture) &
Week 8 &
Equipment Lead &
If not ready, \textbf{RUCK} tests are not admissible; reslot the \textbf{RUCK} gate to next protected window; default \textbf{HOLD}. \\

\textbf{DEP-1C-011} &
7.17 Aquatic Training and Water Safety Equipment — Tier 1 (pool unit verification feasibility; safety readiness) &
Week 8 &
Coach Lead &
If not ready, \textbf{SWIM} gate attempt is not admissible; reslot \textbf{SWIM} to next protected window; do not substitute with non-swim evidence for gate credit. \\

\textbf{DEP-1C-017} &
7.12 Recovery Mobility and Tissue-Care Implements — Tier 1 (ability to execute conservative recovery posture and protect durability) &
Week 1 &
Trainee &
If not available, downshift exposures; preserve validity; if durability degrades, \textbf{HOLD} gates and reslot rather than forcing tests. \\

\end{longtblr}

\endgroup

\subsection*{14. Validity Requirements}
\phantomsection

\begin{itemize}
\item Evidence admissibility is governed by Stage 1.F and Stage 2 validity conditions.
\item Tests must be tagged \textbf{VALID} or \textbf{INVALID} only; \textbf{INVALID} provides no gate credit and requires reslot per Stage 3.E.
\item \textbf{CAL} tests require compliant artifact evidence for \textbf{VALID} gate use per protocol posture.
\item Route/pool/machine identity and relevant settings must be logged for admissibility; missing required context defaults \textbf{INVALID} for gate use.
\item Tool/environment drift near a protected window is treated as comparability risk; default posture is \textbf{HOLD} and reslot until stabilized.
\end{itemize}

\end{landscape}

\subsection*{15. Change Control Notes}
\phantomsection

\begin{itemize}
\item Allowed without patch:
  \begin{itemize}
  \item Reslotting a missed/invalid test to the next protected window within the phase, consistent with Stage 3.E, without stacking additional major events into one window.
  \item Swapping sequencing inside a protected week to preserve spacing intent without changing the test set.
  \end{itemize}
\item Requires patch (Stage 1.F discipline):
  \begin{itemize}
  \item Changing the test set, changing gate focus, adding a test week, extending phase duration, or altering work:deload posture.
  \end{itemize}
\item Missed or invalid tests are corrected by reslot, not by compression or stacking.
\end{itemize}

\subsection*{16. Compliance Checklist}
\phantomsection

\begin{itemize}
\item Pass: All gate tests listed are Stage 2-defined (C\# and \textbf{MET-2B}-\#\#\#) and occur only in protected windows.
\item Pass: Deload/test placement is stated (Week 1 baseline; Weeks 4/8/12 protected windows) and aligned to Stage 3 posture.
\item Pass: Weekly structure table includes Cardio cell structure (Modality; Minutes; RPE + Talk Test; Continuity continuous planned) and does not include minimums in headers.
\item Pass: Constraints include six sessions/week, required element order, Cardio minimum in body text, and 10,000 steps/day in body text.
\item Pass: Dependencies use Stage 7 category IDs and allowed owners.
\item Fail: Weekly Structure table still uses "Minutes: logged" placeholders in gate-relevant structure cells
\end{itemize}

\end{document}